\begin{figure}[H]
    \centering
\begin{tikzpicture}[>=Stealth, arraynode/.style={draw, minimum width=7mm, minimum height=6mm, outer sep=0, inner sep=0}]

\node[circle, draw, fill=lightorange] (x) {$x$};

\node[diamond, draw, right=8mm of x, aspect=2, fill=lightyellow, align=center] (Mx)
    {Learned Model\\$M(x) > \tau$?};

\node[rectangle, draw, below=8mm of Mx, fill=green!15, align=center] (Mxtrue)
    {$x \in S$};

\node[rectangle, draw, right=8mm of Mx, fill=lightyellow, align=center] (fn)
    {False Negative\\Backup Bloom Filter};

\node[rectangle, draw, right=8mm of fn, fill=red!20, align=center] (fnfalse)
    {$x \notin S$};

\node[rectangle, draw, fill=green!15, align=center] (fntrue) at (fn |- Mxtrue)
    {$x \in S$};

% Edges
\draw[->] (x) -- (Mx);
\draw[->] (Mx) -- node[right, font=\scriptsize] {Yes} (Mxtrue);
\draw[->] (Mx) -- node[above, font=\scriptsize] {No}  (fn);
\draw[->] (fn) -- node[above, font=\scriptsize] {No}  (fnfalse);
\draw[->] (fn) -- node[right, font=\scriptsize] {Yes} (fntrue);

\end{tikzpicture}
\caption{Learned Bloom filter. Learned model $M$ is queried with key $x$. If $M$ predicts membership with score greater than $\tau$, output that $x$ is
in the set. Otherwise, output the result of querying a Bloom filter containing the false negatives of $M$. False positives are possible but not false 
negatives.}
\label{fig:csc592_mlsec_lbf}
\end{figure}